\subsection{Acquisition, observation, and the resolution invariant}
\label{sec:acquisition-observation-comparison}
The two classifications distinguish the coordinates that can be acquired
from the metric in which those coordinates are observed. In the monomial
case let $z_a(h)=(\E[g_j^a\mid h])_j$ be the normalized complete Newton
flag, and put $D_a=\operatorname{diag}(d_j)$. Equation~\eqref{eq:leja-physical}
and the fixed spanning query menu give the chain
\[
 h\longmapsto z_a(h)\longmapsto D_a z_a(h)
       \longmapsto\hbox{conditional query probabilities}.
\]
The last map is uniformly conditioned on its active coordinates. Each
initial flag of length at most $p_{n,m}$ is attained with a uniform
positive-history minorization, including at confluent parameter values.
The $d_j$ belong to observation: a zero pivot makes a normalized
coordinate physically invisible, without invalidating its confluent
extension. It is essential here to use a complete flag, rather than
individual divided differences selected after the collision.

In the circular case the corresponding chain is
\[
 h\longmapsto(S_j(h))_j
   \longmapsto(c_j(h))_j=(\tau^jS_j(h))_j
   \longmapsto(Y_j(h))_j=(\tau^{2j}S_j(h))_j.
\]
The normalized $S$-flags have the uniform acquisition minorization;
$c$ consists of the posterior Fourier coefficients; the query metric is
uniformly conditioned in $Y$. Each complex coordinate contributes two
real directions. Thus acquisition and observation contribute separate
factors of $\tau^j$. At zero contrast the normalized extension remains
an auxiliary device and every physical $Y_j$ is zero. These statements
are proved in Lemmas~\ref{lem:circle-attainment} and
\ref{lem:circle-query-metric}, not inferred from Fourier closure.

For nonincreasing real side orders $s_1,\ldots,s_p$, write
\[
 \mathcal Q_M(s;p)=\max_{1\le\ell\le p}
       \left(\frac{\prod_{i=1}^{\ell}s_i}{M}\right)^{2/\ell}.
\]
The monomial law is $\mathcal Q_M(d;p_{n,m})$, since the initial
Leja products are comparable to $\mathcal V_{m,\ell}$. The circular
law is the same expression on the paired list
$\tau^2,\tau^2,\ldots,\tau^{2k},\tau^{2k}$, $k=\min(n,m)$;
its even products give $\mathcal H_{n,m}$ and its odd products add
no larger term. This is a comparison of the two proved laws, not an
assertion that an arbitrary positive detector admits either normal form.
The product-tangent criterion supplies a sufficient attainment test;
the identification of its tangent and its observation scales is a
separate obligation in each experiment.

This identification is the structural content of the classification.
Exact dimension records the number of nonzero directions. The full
profile records their finite-resolution effect jointly in the label
budget and the degenerating parameter, without first fixing a collision
stratum. Its truncation records the restriction imposed by a finite
past. The global covering and causal statements show that these are
operational scales, rather than only local differential invariants.
The separate ambiguity theorem tests the same observation flag against
uncertainty of the prior, and consequently has no past-dimension
truncation. No claim about a sharp memory law follows from a dimension
count, a posterior recurrence, or bounded duality alone.

\subsection{Relation to Bayesian Fourier filtering}
\label{sec:fourier-filtering-comparison}
Exact Fourier representations of Bayesian phase posteriors and their
coefficient updates are established methods. Van den Berg
\cite[Section 2.2, equations (12)--(15)]{VanDenBerg2021} treats
trigonometric posterior representations, their exact updates, and the
complexity caused by coefficient growth. De Neeve, Lebedev, Negnevitsky
and Home \cite{DeNeeveEtAl2025} use reduced-contrast likelihoods and
Fourier-coefficient recursions; see Section II, equations (1)--(3), of
the cited preprint. Their numerical representations in Section IV
address coefficient storage and computation in adaptive phase
estimation. Fourier closure, neighboring-coefficient updates, and the
use of a finite Fourier state are therefore not contributions claimed
here.

The comparison with Theorem~\ref{thm:circle-resolution} concerns a
different experiment, loss, and resource. Our four-cell lookup interface
admits only the specified finite detector commands; every acquisition
trial and failure is counted. A finite menu of physical future queries
is revealed only after the persistent label has been formed. The loss
is squared error in their conditional probabilities, and the resource
is the cardinality $M$ of the persistent state, not the number of
stored real coefficients, the evolution time, or the accuracy of a
phase estimator. A vector of finitely many exact real coefficients
still takes continuously many values and does not by itself meet a
finite-label constraint.

The additional quantitative conclusion is a uniform lower measure on
the normalized acquisition flags, their identification with paired
physical scales $\tau^{2j}$, and matching unconditional and worst-history
finite-label laws. The weighted update then realizes the maximum
checkpoint law with one stage-compatible transducer, without access
to a discarded prefix. These assertions are proved for each fixed
$N$, uniformly in $M$ and $0\le\tau\le\tau_*(N,\eta,\rho)$.
The contrast and model are known to the read-only program; its length
and transient arithmetic are not counted as persistent labels.
Neither the interval nor the comparison constants are claimed uniform
as $N\to\infty$. The cited representation and update formulas are
not identified with this minimax conclusion, nor is our theorem claimed
to subsume the cited phase-estimation results. The circular causal
construction is intrinsic; the effective compiler proved in the
appendices concerns the monomial model.
